% Options for packages loaded elsewhere
\PassOptionsToPackage{unicode}{hyperref}
\PassOptionsToPackage{hyphens}{url}
%
\documentclass[
  ignorenonframetext,
]{beamer}
\usepackage{pgfpages}
\setbeamertemplate{caption}[numbered]
\setbeamertemplate{caption label separator}{: }
\setbeamercolor{caption name}{fg=normal text.fg}
\beamertemplatenavigationsymbolsempty
% Prevent slide breaks in the middle of a paragraph
\widowpenalties 1 10000
\raggedbottom
\setbeamertemplate{part page}{
  \centering
  \begin{beamercolorbox}[sep=16pt,center]{part title}
    \usebeamerfont{part title}\insertpart\par
  \end{beamercolorbox}
}
\setbeamertemplate{section page}{
  \centering
  \begin{beamercolorbox}[sep=12pt,center]{part title}
    \usebeamerfont{section title}\insertsection\par
  \end{beamercolorbox}
}
\setbeamertemplate{subsection page}{
  \centering
  \begin{beamercolorbox}[sep=8pt,center]{part title}
    \usebeamerfont{subsection title}\insertsubsection\par
  \end{beamercolorbox}
}
\AtBeginPart{
  \frame{\partpage}
}
\AtBeginSection{
  \ifbibliography
  \else
    \frame{\sectionpage}
  \fi
}
\AtBeginSubsection{
  \frame{\subsectionpage}
}
\usepackage{amsmath,amssymb}
\usepackage{iftex}
\ifPDFTeX
  \usepackage[T1]{fontenc}
  \usepackage[utf8]{inputenc}
  \usepackage{textcomp} % provide euro and other symbols
\else % if luatex or xetex
  \usepackage{unicode-math} % this also loads fontspec
  \defaultfontfeatures{Scale=MatchLowercase}
  \defaultfontfeatures[\rmfamily]{Ligatures=TeX,Scale=1}
\fi
\usepackage{lmodern}
\ifPDFTeX\else
  % xetex/luatex font selection
\fi
% Use upquote if available, for straight quotes in verbatim environments
\IfFileExists{upquote.sty}{\usepackage{upquote}}{}
\IfFileExists{microtype.sty}{% use microtype if available
  \usepackage[]{microtype}
  \UseMicrotypeSet[protrusion]{basicmath} % disable protrusion for tt fonts
}{}
\makeatletter
\@ifundefined{KOMAClassName}{% if non-KOMA class
  \IfFileExists{parskip.sty}{%
    \usepackage{parskip}
  }{% else
    \setlength{\parindent}{0pt}
    \setlength{\parskip}{6pt plus 2pt minus 1pt}}
}{% if KOMA class
  \KOMAoptions{parskip=half}}
\makeatother
\usepackage{xcolor}
\newif\ifbibliography
\usepackage{longtable,booktabs,array}
\usepackage{calc} % for calculating minipage widths
\usepackage{caption}
% Make caption package work with longtable
\makeatletter
\def\fnum@table{\tablename~\thetable}
\makeatother
\usepackage{graphicx}
\makeatletter
\def\maxwidth{\ifdim\Gin@nat@width>\linewidth\linewidth\else\Gin@nat@width\fi}
\def\maxheight{\ifdim\Gin@nat@height>\textheight\textheight\else\Gin@nat@height\fi}
\makeatother
% Scale images if necessary, so that they will not overflow the page
% margins by default, and it is still possible to overwrite the defaults
% using explicit options in \includegraphics[width, height, ...]{}
\setkeys{Gin}{width=\maxwidth,height=\maxheight,keepaspectratio}
% Set default figure placement to htbp
\makeatletter
\def\fps@figure{htbp}
\makeatother
\setlength{\emergencystretch}{3em} % prevent overfull lines
\providecommand{\tightlist}{%
  \setlength{\itemsep}{0pt}\setlength{\parskip}{0pt}}
\setcounter{secnumdepth}{-\maxdimen} % remove section numbering
\usepackage{quiver}
\usepackage{mathpartir}
% \usepackage{enumitem}
\usepackage{wrapfig}
\usepackage{fancyvrb}
\usepackage{comment}
\usepackage{array}

\usepackage{stmaryrd}
\usepackage{graphicx}
\newcommand{\To}{\Rightarrow}
\newcommand{\inl}{\mathsf{inl}}
\newcommand{\inr}{\mathsf{inr}}
\newcommand{\altbar}{\mathrel{\bf \,\mid\,}}

\newcommand{\ob}{\text{Ob}}


\newcommand{\extlc}{\text{Ext-}\lambda}
\newcommand{\extlcm}{\text{Ext-}\lambda^{-\text{trans}}}
\newcommand{\extlcmm}{\text{Ext-}\lambda^{-\text{trans}-\text{cast}}}
\newcommand{\extlcprime}{\text{Ext-}\lambda'}
\newcommand{\intlc}{\text{Int-}\lambda}
\newcommand{\intlcbisim}{\text{Int$_\approx$-}\lambda}
\newcommand{\erase}[1]{\lfloor {#1} \rfloor}


\newcommand{\uarrowl}{\mathrel{\rotatebox[origin=c]{-30}{$\leftarrowtail$}}}
\newcommand{\uarrowr}{\mathrel{\rotatebox[origin=c]{60}{$\leftarrowtail$}}}
\newcommand{\darrowl}{\mathrel{\rotatebox[origin=c]{30}{$\twoheadleftarrow$}}}
\newcommand{\darrowr}{\mathrel{\rotatebox[origin=c]{120}{$\twoheadleftarrow$}}}
\newcommand{\vuarrow}{\mathrel{\rotatebox[origin=c]{-90}{$\leftarrowtail$}}}
\newcommand{\vdarrow}{\mathrel{\rotatebox[origin=c]{90}{$\twoheadleftarrow$}}}

% Types, terms, and precision 

\newcommand{\dyn}{{?}}
\newcommand{\nat}{\text{Nat}}
\newcommand{\bool}{\text{Bool}}
\newcommand{\ra}{\rightharpoonup}
\newcommand{\Ret}[1]{\mathsf{Ret\,}{#1}}
\newcommand{\hole}[1]{\bullet \colon {#1}}
\newcommand{\dyntodyn}{\dyn \ra\, \dyn}


\newcommand{\up}[2]{\langle{#2}\uarrowl{#1}\rangle}
\newcommand{\dn}[2]{\langle{#1}\darrowl{#2}\rangle}

\newcommand{\upc}[1]{\text{up}\,{#1}\,}
\newcommand{\dnc}[1]{\text{dn}\,{#1}\,}


% \newcommand{\ret}{\mathsf{ret}}
\newcommand{\err}{\mho}
\newcommand{\zro}{\textsf{zro}}
\newcommand{\suc}{\textsf{suc}}
\newcommand{\lda}[2]{\lambda {#1} . {#2}}

\newcommand{\injarr}[1]{\textsf{Inj}_\ra ({#1})}
\newcommand{\injnat}[1]{\textsf{Inj}_\text{nat} ({#1})}
\newcommand{\casenat}[4]{\text{Case}_\text{nat} ({#1}) \{ \text{no} \to {#2} \altbar \text{nat}({#3}) \to {#4} \}}
\newcommand{\casearr}[4]{\text{Case}_\ra ({#1}) \{ \text{no} \to {#2} \altbar \text{fun}({#3}) \to {#4} \}}
\newcommand{\casedyn}[5]{\text{Case}_\Dyn ({#1}) \{ \text{nat}({#2}) \to {#3} \altbar \text{fun}({#4}) \to {#5} \}}
%\newcommand{\bind}[3]{\text{var } {#1} = {#2} \text{ in } {#3}}
\newcommand{\matchnat}[4]{\text{match } {#1} \text{ with } \{ \textsf {zero} \Rightarrow {#2} \altbar \textsf{ suc } {#3} \Rightarrow {#4} \}}

\newcommand{\refl}{\text{refl}}

\newcommand{\Lift}{\text{Lift}}

%\newcommand{\rel}{\circ\hspace{-4px}-\hspace{-4px}\bullet}
%% \newcommand{\rel}{\mathrel{\circ\mkern-6mu-\mkern-6mu\bullet}}
%% \def\rel{\mathrel{\joinrel=}}
\newcommand{\rel}{\mathrel{{\relbar}\hspace{-2px}{\relbar}\hspace{-2px}{\shortmid}\hspace{-2px}{\relbar}\hspace{-2px}{\relbar}}}
% \newcommand{\wand}{\mathrel{-\mkern-6mu*}}
\newcommand{\wand}{\mathrel{\multimap}}
\newcommand{\ltdyn}{\sqsubseteq}
\newcommand{\gtdyn}{\sqsupseteq}
\newcommand{\equidyn}{\mathrel{\gtdyn\ltdyn}}
\newcommand{\gamlt}{\Gamma^\ltdyn}
\newcommand{\deltalt}{\Delta^\ltdyn}
\newcommand{\relcomp}{\odot}

\newcommand{\hasty}[3]{{#1} \vdash {#2} \colon {#3}}
\newcommand{\vhasty}[3]{{#1} \vdash^v {#2} \colon {#3}}
\newcommand{\phasty}[3]{{#1} \vdash^c {#2} \colon {#3}}
\newcommand{\etmprec}[4]{{#1} \vdash {#2} \ltdyn_e {#3} \colon {#4}}
\newcommand{\itmprec}[4]{{#1} \vdash {#2} \ltdyn_i {#3} \colon {#4}}
\newcommand{\etmequidyn}[4]{{#1} \vdash {#2} \equidyn_e {#3} \colon {#4}}
\newcommand{\itmequidyn}[4]{{#1} \vdash {#2} \equidyn_i {#3} \colon {#4}}
\newcommand{\synbisim}{\approx_{\text{syn}}}

\newcommand{\Dwn}{\Downarrow}
\newcommand{\qte}[1]{\text{quote}({#1})}

\newcommand{\elab}[1]{\text{Elab}({#1})}

% Perturbations
\newcommand{\pertp}{\text{Pert}^\text{P}}
\newcommand{\perte}{\text{Pert}^\text{E}}
\newcommand{\pertdyn}[2]{\text{pert-dyn}({#1}, {#2})}
\newcommand{\delaypert}[1]{\text{delay-pert}({#1})}

\newcommand{\pertc}{\text{Pert}_{\text{C}}}
\newcommand{\pertv}{\text{Pert}_{\text{V}}}


% SGDT and Intensional Stuff

\newcommand{\later}{{\vartriangleright}}
\newcommand{\laterhs}{{\later}}
\newcommand{\type}{\texttt{Type}}
\newcommand{\lob}{\text{L\"{o}b}}
\newcommand{\tick}{\mathsf{tick}}
\newcommand{\nxt}{\mathsf{next}}
\newcommand{\fix}{\mathsf{fix}}
\newcommand{\kpa}{\kappa}

% Model-related stuff
\newcommand{\calV}{\mathcal{V}}
\newcommand{\calE}{\mathcal{E}}
\newcommand{\calVk}{\mathcal{V}_k}
\newcommand{\calEk}{\mathcal{E}_k}
\newcommand{\tok}{\mathrel{\mathop{\to}\limits^{\textrm{\footnotesize k}}}}
\newcommand{\timesk}{\mathrel{\mathop{\times}\limits^{\textrm{k}}}}
\newcommand{\Fk}{\mathrel{\mathop{F}\limits^{\textrm{k}}}}
\newcommand{\Uk}{\mathrel{\mathop{U}\limits^{\textrm{k}}}}

\newcommand{\op}[1]{{#1}^{\textrm{op}}}
\newcommand{\calC}{\mathcal{C}}
\newcommand{\Set}{\mathsf{Set}}
\newcommand{\ErrDom}{\mathsf{ErrDom}}
\newcommand{\PreDom}{\mathsf{PreDom}}
\newcommand{\Yo}{\mathsf{Yo}}
\newcommand{\Hom}{\mathsf{Hom}}
\newcommand{\calS}{\mathcal{S}}
\newcommand{\gfix}{\texttt{gfix}}
\newcommand{\qfix}{\texttt{qfix}}
\newcommand{\calU}{\mathcal{U}}
\newcommand{\laterhat}{\widehat{\later}}
\newcommand{\El}{\mathsf{El}}
\newcommand{\Clock}{\mathsf{Clock}}

\newcommand{\Machine}[1]{\mathsf{Machine}\, {#1}}


% Predomains and EP pairs
\newcommand{\Nat}{\mathsf{Nat}}
\newcommand{\Dyn}{\mathsf{Dyn}}
\newcommand{\ty}[1]{\langle {#1} \rangle}
\newcommand{\li}{L_\mho}
\newcommand{\liclk}[1]{L_\mho [{#1}]}

\newcommand{\ext}[2]{\text{ext}\,{#1}\,{#2}}
\newcommand{\map}[2]{\text{map}\,{#1}\,{#2}}

\newcommand{\ltls}{\ltdyn}
\newcommand{\bisim}{\approx}
\newcommand{\semlt}{\le}
\newcommand{\semltbad}{\lesssim}

%\newcommand{\injarr}{\textsf{Inj}_\to}
%\newcommand{\injnat}{\textsf{Inj}_\mathbb{N}}

\newcommand{\id}{\mathsf{id}}
\newcommand{\ep}{\leadsto}

\newcommand{\emb}[2]{\mathsf{emb}_{#1}({#2})}
\newcommand{\proj}[2]{\mathsf{proj}_{#1}({#2})}

\newcommand{\monto}{\to_m}


% Notation for wait functions
\newcommand{\wre}{w_r^e}
\newcommand{\wle}{w_l^e}
\newcommand{\wrp}{w_r^p}
\newcommand{\wlp}{w_l^p}

\newcommand{\sem}[1]{\llbracket {#1} \rrbracket}
\newcommand{\semgl}[1]{{\sem{#1}}^\text{gl}}


% Denotational model
\newcommand{\errdom}{\textsf{ErrDom}}

\newcommand{\ptb}{\text{ptb}}
\newcommand{\push}{\text{push}}
\newcommand{\pull}{\text{pull}}

\newcommand{\pv}{P^{\mathcal{V}}}
\newcommand{\pe}{P^{\mathcal{E}}}
\newcommand{\ptbv}{\text{ptb}^\mathcal{V}}
\newcommand{\ptbe}{\text{ptb}^\mathcal{E}}

\newcommand{\Ppure}{P}
\newcommand{\Pk}{P^K}
\newcommand{\ptbk}{\text{ptb}^K}


\newcommand{\upf}{\text{up}}
\newcommand{\dnf}{\text{dn}
}
\newcommand{\arr}{\to}
\newcommand{\comp}{}

\newcommand{\ltsq}[2]{\mathrel{\ltdyn^{#1}_{#2}}}
\newcommand{\bisimsq}[2]{\mathrel{\bisim^{#1}_{#2}}}
\newcommand{\ltsqbisim}[2]{\mathrel{{\widetilde{\ltdyn}}^{#1}_{#2}}}
\newcommand{\ltbisim}{\mathrel{\widetilde{\ltdyn}}}
\newcommand{\vf}{\mathcal{V}_f}
\newcommand{\vsq}{\mathcal{V}_{sq}}
\newcommand{\ef}{\mathcal{E}_f}
\newcommand{\esq}{\mathcal{E}_{sq}}

\newcommand{\morbisimid}[1]{\text{Endo}_\bisim({#1})}


\newcommand{\sv}{s_{\mathcal{V}}}
\newcommand{\tv}{t_{\mathcal{V}}}
\newcommand{\rv}{r_{\mathcal{V}}}
\newcommand{\se}{s_{\mathcal{E}}}
\newcommand{\te}{t_{\mathcal{E}}}
\newcommand{\re}{r_{\mathcal{E}}}

\newcommand{\Ff}{F_f}
\newcommand{\Fsq}{F_{sq}}
\newcommand{\Uf}{U_f}
\newcommand{\Usq}{U_{sq}}
\newcommand{\arrf}{\arr_f}
\newcommand{\arrsq}{\arr_{sq}}

\newcommand{\vsim}{\mathcal{V}_{\bisim}}
\newcommand{\esim}{\mathcal{E}_{\bisim}}
\newcommand{\svsim}{s_{\mathcal{V}}^\bisim}
\newcommand{\tvsim}{t_{\mathcal{V}}^\bisim}
\newcommand{\rvsim}{r_{\mathcal{V}}^\bisim}
\newcommand{\sesim}{s_{\mathcal{E}}^\bisim}
\newcommand{\tesim}{t_{\mathcal{E}}^\bisim}
\newcommand{\resim}{r_{\mathcal{E}}^\bisim}



\newcommand{\ve}{\mathcal{V}_e}
\newcommand{\ee}{\mathcal{E}_e}

\newcommand{\vr}{\mathcal{V}_r}
\newcommand{\er}{\mathcal{E}_r}

\newcommand{\relatedin}[3]{{#1}\, {#3}\, {#2}}
\newcommand{\binrel}[1]{\mathbin{#1}}


\newcommand{\da}{\downarrow}

\newcommand{\ret}{\text{ret}}
\newcommand{\bind}[3]{{#1} <- {#2}\, ; \, {#3}}

\newcommand{\piv}{\Pi^\mathcal{V}}
\newcommand{\pie}{\Pi^\mathcal{E}}

\newcommand{\upl}{\textsc{UpL}}
\newcommand{\upr}{\textsc{UpR}}
\newcommand{\dnl}{\textsc{DnL}}
\newcommand{\dnr}{\textsc{DnR}}

\newcommand{\delre}{\delta^{r,e}}
\newcommand{\delle}{\delta^{l,e}}
\newcommand{\delrp}{\delta^{r,p}}
\newcommand{\dellp}{\delta^{l,p}}

\newcommand{\qordeq}{\bisim}

\newcommand{\inat}{\text{Inj}_\mathbb{N}}
\newcommand{\itimes}{\text{Inj}_\times}
\newcommand{\iarr}{\text{Inj}_\to}

\newcommand{\tnat}{\mathsf{nat}}
\newcommand{\ttimes}{\mathsf{times}}
\newcommand{\tfun}{\mathsf{fun}}

\newcommand{\cased}[7]{
    \begin{align*}
    \text{Case}_D ({#1}) &\text{ of } \{ \\
        &\altbar \text{nat}({#2}) \to {#3}  \\
        &\altbar \text{times}({#4}) \to {#5}  \\
        &\altbar \text{fun}({#6}) \to {#7} \}
    \end{align*}
}

\newcommand{\inone}{\mathsf{in}_1}
\newcommand{\intwo}{\mathsf{in}_2}
\newcommand{\inthree}{\mathsf{in}_3}
\newcommand{\infour}{\mathsf{in}_4}
\newcommand{\infive}{\mathsf{in}_5}


\newcommand{\delay}{\text{Delay}}
\newcommand{\ledelay}{\le^{\text{Del}}}
\newcommand{\bisimdelay}{\bisim^{\text{Del}}}
\newcommand{\tnow}{\mathsf{now}}
\newcommand{\tlater}{\mathsf{later}}
\newcommand{\Da}{\Downarrow}
\ifLuaTeX
  \usepackage{selnolig}  % disable illegal ligatures
\fi
\IfFileExists{bookmark.sty}{\usepackage{bookmark}}{\usepackage{hyperref}}
\IfFileExists{xurl.sty}{\usepackage{xurl}}{} % add URL line breaks if available
\urlstyle{same}
\hypersetup{
  pdftitle={Denotational Semantics of Gradual Typing using Synthetic Guarded Domain Theory},
  pdfauthor={Eric Giovannini, Tingting Ding, and Max S. New},
  hidelinks,
  pdfcreator={LaTeX via pandoc}}

\title{Denotational Semantics of Gradual Typing using Synthetic Guarded
Domain Theory}
\author{Eric Giovannini, Tingting Ding, and Max S. New}
\date{November 2024}

\begin{document}
\frame{\titlepage}

\hypertarget{introduction-and-background}{%
\section{Introduction and
Background}\label{introduction-and-background}}

\begin{frame}{Big Idea}
\protect\hypertarget{big-idea}{}
This work is about constructing a \textbf{denotational semantics} for
\textbf{gradually-typed programming languages}.
\end{frame}

\begin{frame}{Gradual Typing}
\protect\hypertarget{gradual-typing}{}
\begin{itemize}
\tightlist
\item
  Idea: Combine static and dynamic typing disciplines in the same
  language
\end{itemize}
\end{frame}

\begin{frame}{Desired Properties of Gradually Typed Languages}
\protect\hypertarget{desired-properties-of-gradually-typed-languages}{}
\begin{itemize}
\tightlist
\item
  \textbf{Graduality}: Migrating from a static to a dynamic typing
  discipline should not change the semantics of the program except to
  catch more errors
\end{itemize}

\pause

\includegraphics{fig1.drawio.png}\\

\pause

\begin{itemize}
\tightlist
\item
  \textbf{Soundness} of the equational theory: justifies program
  optimizations
\end{itemize}
\end{frame}

\begin{frame}{Our Goal}
\protect\hypertarget{our-goal}{}
\textbf{Provide an \emph{expressive}, \emph{reusable},
\emph{compositional} semantic framework for defining well-behaved
semantics of gradually typed programs.}
\end{frame}

\begin{frame}{Denotational Semantics of Gradually-Typed Languages}
\protect\hypertarget{denotational-semantics-of-gradually-typed-languages}{}
\begin{itemize}
\item
  Even simple gradually-typed languages are nontrivial to model!

  \begin{itemize}
  \tightlist
  \item
    Recursion and recursive types through the presence of dynamic types
    \(\rightarrow\) \textbf{domain theory}
  \end{itemize}

  \pause

  \begin{itemize}
  \tightlist
  \item
    Effects in the form of divergence and errors \(\rightarrow\)
    \textbf{monads} or \textbf{adjunctions}
  \end{itemize}

  \pause

  \begin{itemize}
  \tightlist
  \item
    Relational models to capture the graduality property \(\rightarrow\)
    \textbf{reflexive graph categories} or \textbf{double categories}
  \end{itemize}
\end{itemize}

\pause

\begin{itemize}
\tightlist
\item
  Prior work: denotational models using classical domain theory (New and
  Licata 2018)

  \begin{itemize}
  \tightlist
  \item
    Benefit: compositional and reusable
  \item
    Limitation: can't scale to advanced language features such as
    higher-order store
  \end{itemize}
\end{itemize}
\end{frame}

\begin{frame}{Guarded Domain Theory}
\protect\hypertarget{guarded-domain-theory}{}
\begin{itemize}
\item
  Alternative to classical domain theory that can scale to arbitrarily
  complex definitions
\item
  Based on an entirely different foundations, most commonly
  ``step-indexed sets'\,', i.e., objects in the \emph{topos of trees}
  \cite{birkedal-mogelberg-schwinghammer-stovring2011}.

  \begin{itemize}
  \tightlist
  \item
    A family \(\{X_n\}_{n \in \mathbb{N}}\) of sets along with
    restriction functions \(r_n : X_{n+1} \to X_n\) for all \(n\).\\
  \item
    Idea: infinite sequence of increasingly precise approximations to
    the true type being modeled.
  \end{itemize}
\end{itemize}
\end{frame}

\begin{frame}{Guarded Domain Theory, Continued}
\protect\hypertarget{guarded-domain-theory-continued}{}
\begin{itemize}
\tightlist
\item
  Operator \(\later\) (pronounced ``later'\,')
\item
  Delays the approximation by one step
\item
  Any guarded domain equation \(X \cong F(\later X)\) has a unique
  solution.
\end{itemize}
\end{frame}

\begin{frame}{Guarded Type Theory}
\protect\hypertarget{guarded-type-theory}{}
\begin{itemize}
\tightlist
\item
  Synthetic presentation of guarded domain theory
\item
  \(\later\) is taken as a primitive operation on types
\item
  Axiom that \emph{guarded} domain equations have a (necessarily unique)
  solution.
\item
  Benefit: we model an object language type as a \emph{set} instead of a
  step-indexed set
\end{itemize}
\end{frame}

\begin{frame}{Adapting the New-Licata Approach}
\protect\hypertarget{adapting-the-new-licata-approach}{}
\begin{itemize}
\tightlist
\item
  \textbf{Motivating question}: can we adapt the denotational model of
  New-Licata to guarded domain theory?
\end{itemize}

\pause

\begin{itemize}
\tightlist
\item
  Key issue: \emph{intensional} nature of the semantics

  \begin{itemize}
  \tightlist
  \item
    Unfolding the dynamic type involves an observable computational step
  \end{itemize}
\end{itemize}

\pause

\begin{itemize}
\item
  Casts take observable steps \(\rightarrow\) graduality must be
  insensitive to steps
\item
  So the model must allow for reasoning \emph{up to weak bisimilarity}

  \begin{itemize}
  \tightlist
  \item
    Two terms are \textbf{weakly bisimilar} if they differ only in their
    number of computational steps
  \end{itemize}
\end{itemize}
\end{frame}

\begin{frame}{Key Issue in Constructing the Model}
\protect\hypertarget{key-issue-in-constructing-the-model}{}
\begin{itemize}
\tightlist
\item
  Some amount of transitive reasoning is necessary to define a
  syntax-independent model of gradual typing
\end{itemize}

\pause

\begin{itemize}
\tightlist
\item
  In the guarded setting, it is not possible to reason transitively up
  to weak bisimilarity! (Follows from a no-go theorem we establish)
\end{itemize}

\pause

\begin{itemize}
\tightlist
\item
  Our solution: combination of strong error ordering + weak bisimilarity
\end{itemize}

\pause

\begin{longtable}[]{@{}cc@{}}
\toprule\noalign{}
Strong error ordering & Weak Bisimilarity \\
\midrule\noalign{}
\endhead
sensitive to steps & insensitive to steps \\
transitive reasoning & \emph{not} transitive \\
\bottomrule\noalign{}
\end{longtable}

\pause

\begin{itemize}
\tightlist
\item
  \textbf{syntactic perturbations}: explicit synchronizations that we
  manipulate syntactically
\end{itemize}
\end{frame}

\begin{frame}{Adequacy}
\protect\hypertarget{adequacy}{}
\end{frame}

\begin{frame}{Contributions}
\protect\hypertarget{contributions}{}
\begin{itemize}
\item
  We define a denotational semantics for a simple gradually-typed cast
  calculus in guarded type theory.
\item
  We prove that the model is adequate for the graduality property.
\item
  Most of our core theorems have been verified in Guarded Cubical Agda.
\end{itemize}
\end{frame}

\begin{frame}{Outline of Remainder of the Talk}
\protect\hypertarget{outline-of-remainder-of-the-talk}{}
\begin{enumerate}
\tightlist
\item
  Describe the syntax of the cast calculus that we will interpret into
  our model.
\item
  Give a model for the \emph{terms} of our cast calculus.
\item
  Extend the model to a relational model.
\item
  Discuss related and future work.
\end{enumerate}
\end{frame}

\hypertarget{syntax}{%
\section{Syntax}\label{syntax}}

\begin{frame}{Types and Terms}
\protect\hypertarget{types-and-terms}{}
\begin{figure}
  \begin{mathpar}
  \begin{array}{rcl}
    \text{Types } A &::=& \nat \altbar \,\dyn \altbar A \ra A' \altbar A \times A'\\
    \text{Type Precision } c &::=& r(A) \altbar c c' \altbar \iarr \\
        &&\altbar \inat \altbar \itimes \altbar c \ra c' \altbar c \times c'\\
    \text{Values } V &::=& x \altbar \upc c V \altbar \zro \altbar \suc\, V \altbar \lda{x}{M} \altbar (V,V') \\ 
    \text{Terms } M,N &::=& \err\altbar \upc c M \altbar \dnc c M \altbar \zro \altbar \suc\, M \altbar \lda{x}{M} \\ 
        &&\altbar M\, N \altbar (M,N) \altbar \textrm{let } (x,y) = M \textrm{ in } N\\
    \text{Contexts } \Gamma &::= &\cdot \altbar \Gamma, x : A \\
    \text{Ctx Precision } \Delta &::=& \cdot\altbar \Delta,x:c
  \end{array}
  \end{mathpar}
  \caption{GTLC Cast Calculus Syntax}
  \label{fig:gtlc-syntax}
\end{figure}
\end{frame}

\begin{frame}{Typing Rules (Selected)}
\protect\hypertarget{typing-rules-selected}{}
\begin{figure}
  \begin{mathpar}
 
  \inferrule
  {\Gamma \vdash M : A \and c : A \ltdyn A'}
  {\Gamma \vdash \upc c M : A'}

  \inferrule
  {\Gamma \vdash N : A' \and c : A \ltdyn A'}
  {\Gamma \vdash \dnc c N : A}

  \inferrule{}{\Gamma \vdash \mho : A}
  
  \end{mathpar}
  \caption{GTLC Typing Rules (Selected)}
  \label{fig:gtlc-typing}
\end{figure}
\end{frame}

\begin{frame}{Type Precision}
\protect\hypertarget{type-precision}{}
\begin{figure}
  \begin{mathpar}
    \inferrule{}{r(A) : A \ltdyn A} \and
    \inferrule{c : A \ltdyn A' \and c' : A' \ltdyn A''}{cc' : A \ltdyn A''} \\
    \inferrule{}{\iarr \colon \dyn \ra \dyn \ltdyn \dyn}\and
    \inferrule{}{\inat \colon \nat \ltdyn \dyn} \and
    \inferrule{}{\itimes \colon \dyn \times \dyn \ltdyn \dyn} \\
    \scalebox{0.80}{\inferrule{c_i : A_i \ltdyn A_i' \and c_o : A_o \ltdyn A_o'}{c_i \ra c_o : (A_i \ra A_o) \ltdyn (A_i' \ra A_o')}}\and
    \scalebox{0.80}{\inferrule{c_1 : A_1 \ltdyn A_1' \and c_2 : A_2 \ltdyn A_2'}{c_1 \times c_2 : (A_1 \times A_2) \ltdyn (A_1' \times A_2')}}
  \end{mathpar}
  \caption{GTLC Type Precision Derivations}
  \label{fig:gtlc-type-prec}
\end{figure}
\end{frame}

\begin{frame}{Type Precision Equivalence}
\protect\hypertarget{type-precision-equivalence}{}
\begin{itemize}
\tightlist
\item
  Equational theory \(c \equiv c'\) that implies that the corresponding
  casts are weakly bisimilar in the semantics
\end{itemize}

\begin{figure}
  \begin{mathpar}
     r(A)c \equiv c\and
     c \equiv cr(A')\and
     c(c'c'') \equiv (cc')c''\and
     r(A_i \ra A_o) \equiv r(A_i) \ra r(A_o)\and
     r(A_1\times A_2) \equiv r(A_1) \times r(A_2)\and
     (c_i \ra c_o)(c_i' \ra c_o')\equiv (c_ic_i' \ra c_oc_o') \and
     (c_1\times c_2)(c_1'\times c_2')\equiv (c_1c_1' \times c_2c_2')
  \end{mathpar}
  \caption{GTLC Type Precision Equivalence}
  \label{fig:gtlc-type-prec-equiv}
\end{figure}
\end{frame}

\begin{frame}{Term Precision}
\protect\hypertarget{term-precision}{}
\begin{itemize}
\item
  Extension of type precision to terms
\item
  \(\Delta \vdash M \ltdyn M' : c\) where \(\Delta\) is a context where
  variables are assigned to type precision derivations.
\end{itemize}

\begin{figure}
  \begin{mathpar}
  %(\lambda x. M)(V) = M[V/x] \and (V : A \ra A') = \lambda x. V\,x\\

  % \textrm{let } (x,y) = (V,V') \textrm{ in } N = N[V/x,V'/y] \and
  % M[V:A\times A'/p] = \textrm{let } (x,y) = V \textrm{ in } M[(x,y)/p]

  \inferrule*[right=EquivTyPrec]
  {\Delta\vdash M \ltdyn M' : c \and c \equiv c'}
  {\Delta\vdash M \ltdyn M' : c'}

  \inferrule*[right=ErrBot]
  {}
  {\Delta \vdash \mho \ltdyn M : c}\\

  \inferrule*[right=UpL]
  {M \ltdyn M' : cc_r}
  {\upc {c} M \ltdyn M' : c_r}

  \inferrule*[right=UpR]
  {M \ltdyn M' : c_l}
  {M \ltdyn \upc {c} M' : c_lc}

  \inferrule*[right=DnL]
  {M \ltdyn M' : c_r}
  {\dnc {c} M \ltdyn M' : cc_r}

  \inferrule*[right=DnR]
  {M \ltdyn M' : c_lc}
  {M \ltdyn \dnc {c} M' : c_l}
  \end{mathpar}
  \caption{Term Precision Rules (Selected)}
  \label{fig:term-prec}
\end{figure}
\end{frame}

\hypertarget{a-simple-term-model-in-guarded-type-theory}{%
\section{A Simple Term Model in Guarded Type
Theory}\label{a-simple-term-model-in-guarded-type-theory}}

\begin{frame}{More about Guarded Type Theory}
\protect\hypertarget{more-about-guarded-type-theory}{}
\begin{itemize}
\tightlist
\item
  Recall: Operator \(\later : \type \to \type\)
\item
  \emph{Guarded} fixpoint operator \(\fix : (\later A \to A) \to A\)
  satisfying \(\fix\, f = f (\nxt (\fix\, f))\)

  \begin{itemize}
  \tightlist
  \item
    The recursive occurrence is guarded by a \(\later\)
  \item
    Constructing a Prop in this manner corresponds to \(\lob\)-induction
  \end{itemize}
\end{itemize}
\end{frame}

\begin{frame}{Ticked Cubical Type Theory}
\protect\hypertarget{ticked-cubical-type-theory}{}
\begin{itemize}
\tightlist
\item
  Abstract sort \(\tick\)
\item
  \(\later A\) is modeled as \(\tick \to A\)
\item
  The type \(A\) is allowed to depend on \(t\); we write \(\later_t A\)
\item
  Rules are similar to those of ordinary \(\Pi\) types
\item
  Restriction on tick application to prevent writing a term
  \(\later \later A \to \later A\)
\item
  Notation: \(M_t\) for tick application
\end{itemize}
\end{frame}

\begin{frame}{A CBPV Model}
\protect\hypertarget{a-cbpv-model}{}
\begin{itemize}
\tightlist
\item
  Goal of this section: define a model for only the \textbf{types and
  terms} of GTLC
\end{itemize}

\pause

\begin{itemize}
\tightlist
\item
  We follow the structure of Levy's Call-by-Push-Value \cite{levy99}

  \begin{itemize}
  \tightlist
  \item
    Refinement of Moggi's monadic semantics \cite{MOGGI199155} that
    decomposes Moggi's monad \(T\) into an adjunction \(T = UF\):
  \end{itemize}
\end{itemize}

% https://q.uiver.app/#q=WzAsNixbMCwwLCJcXGNhbFYiXSxbMiwwLCJcXGNhbEUiXSxbMCwxLCJcXHRleHR7VmFsdWUgVHlwZXN9Il0sWzAsMiwiXFx0ZXh0e1B1cmUgbW9ycGhpc21zfSJdLFsyLDEsIlxcdGV4dHtDb21wdXRhdGlvbiB0eXBlc30iXSxbMiwyLCJcXHRleHR7SG9tb21vcnBoaXNtc30iXSxbMCwxLCJGIiwwLHsiY3VydmUiOi0zfV0sWzEsMCwiVSIsMCx7ImN1cnZlIjotM31dLFs2LDcsIiIsMCx7ImxldmVsIjoxLCJzdHlsZSI6eyJuYW1lIjoiYWRqdW5jdGlvbiJ9fV1d
\[\begin{tikzcd}[ampersand replacement=\&]
    \calV \&\& \calE \\
    {\text{Value Types}} \&\& {\text{Computation types}} \\
    {\text{Pure morphisms}} \&\& {\text{Homomorphisms}}
    \arrow[""{name=0, anchor=center, inner sep=0}, "F", curve={height=-18pt}, from=1-1, to=1-3]
    \arrow[""{name=1, anchor=center, inner sep=0}, "U", curve={height=-18pt}, from=1-3, to=1-1]
    \arrow["\dashv"{anchor=center, rotate=-90}, draw=none, from=0, to=1]
\end{tikzcd}\]

\pause

\begin{itemize}
\tightlist
\item
  \(A \rightharpoonup A'\) decomposes into \(U(A \to F A')\) where
  \(\arr : \calV^{op} \times \calE \to \calE\)
\end{itemize}
\end{frame}

\begin{frame}{Value and Computation Categories}
\protect\hypertarget{value-and-computation-categories}{}
\begin{itemize}
\tightlist
\item
  Value category: sets and functions
\item
  Computation category: sets equipped with algebraic structure, and
  homomorphisms that preserve this structure
\item
  Effects we model: errors and computational steps
\end{itemize}

\pause

\begin{definition}[Simple Error Domains]
  A (simple) error domain $B$ consists of
  \begin{enumerate}
  \item A carrier set $UB$
  \item An element $\mho_{B} : UB$ representing error
  \item A function $\theta_B : \laterhs UB \to UB$ modeling a computational step
  \end{enumerate}
\end{definition}
\end{frame}

\begin{frame}{The Free Error Domain}
\protect\hypertarget{the-free-error-domain}{}
\begin{definition}[Free error domain]
  For a set $A$, we define the (carrier of) \emph{free error domain} $U(\li A)$ as the unique solution to the guarded domain equation:
  \[ U(\li A) \cong A + 1\, + \laterhs U(\li A). \]
  We use the following notation for the three constructors:
  \begin{enumerate}
  \item $\eta \colon A \to U(\li A)$
  \item $\mho \colon U(\li A)$
  \item $\theta \colon \laterhs U(\li A) \to U(\li A)$
  \end{enumerate}
\end{definition}

\pause

\begin{itemize}
\tightlist
\item
  \(\mho\) and \(\theta\) provide the error domain structure for
  \(\li A\)
\end{itemize}

\pause

\begin{itemize}
\tightlist
\item
  Satisfies the following universal property of being free: any function
  \(f : U(\li A) \to UB\) uniquely extends to a homomorphism
  \(f^\dagger : \li A \multimap B\) satisfying
  \(Uf^\dagger \circ \eta = f\).
\end{itemize}

\pause

\begin{itemize}
\tightlist
\item
  We will write \(\theta_t(\dots)\) to mean
  \(\theta (\lambda t. \dots)\).
\end{itemize}
\end{frame}

\begin{frame}{Modeling the Type Structure}
\protect\hypertarget{modeling-the-type-structure}{}
\begin{itemize}
\tightlist
\item
  \(\nat\) denotes the set of natural numbers
\item
  \(\times\) denotes the cartesian product of sets.
\item
  CBV function type \(A \rightharpoonup A'\) denotes the CBPV
  decomposition \(U(\sem{A} \to \li \sem{A'})\)

  \begin{itemize}
  \tightlist
  \item
    \(A \to B\) is the functions from \(A\) to \(UB\) with the obvious
    point-wise algebraic structure
  \end{itemize}
\end{itemize}
\end{frame}

\begin{frame}{The Dynamic Type}
\protect\hypertarget{the-dynamic-type}{}
\textbf{Classical domain theory}:

\[D \cong \Nat + (D \times D) + U(D \to \li D)\]

\begin{itemize}
\tightlist
\item
  No inductive or coinductive solutions
\end{itemize}

\pause

\textbf{Guarded domain theory}:

\begin{equation}\label{eq:dyn}
D \cong \Nat + (D \times D)\, + \laterhs U(D \to \li D).
\end{equation}

\pause

\begin{itemize}
\tightlist
\item
  We solve this equation by a mixture of inductive types and guarded
  fixed points.
\end{itemize}
\end{frame}

\begin{frame}{Solving the Equation for the Dynamic Type}
\protect\hypertarget{solving-the-equation-for-the-dynamic-type}{}
\begin{itemize}
\tightlist
\item
  Consider the parameterized inductive type
\end{itemize}

\[ D'\, X := \mu T. \Nat + (T \times T) + X.\]

\pause

\begin{itemize}
\tightlist
\item
  We construct \(D\) as the unique solution to the guarded domain
  equation
\end{itemize}

\[D \cong D'(\laterhs U(D \to \li D))\]

\pause

\begin{itemize}
\tightlist
\item
  Expanding the guarded fixed point and least fixed point property gives
  us that \(D\) satisfies Equation \(\ref{eq:dyn}\).
\end{itemize}

\pause

\begin{itemize}
\tightlist
\item
  We refer to the three injections for numbers, pairs and functions as
  \(\inat, \itimes, \iarr\) respectively.
\end{itemize}
\end{frame}

\begin{frame}{Interpreting the Terms}
\protect\hypertarget{interpreting-the-terms}{}
\begin{itemize}
\item
  Terms \(x:A_1,\ldots \vdash M : A\) are effectful functions
  \(\sem{M}: \sem{A_1}\times \cdots \to U\li\sem{A}\)
\item
  Values \(x:A_1,\ldots \vdash V : A\) are pure functions
  \(\sem{V} : \sem{A_1}\times\cdots \to \sem{A}\).
\item
  Only nontrivial terms: \textbf{casts}
\end{itemize}
\end{frame}

\begin{frame}{Interpreting Casts}
\protect\hypertarget{interpreting-casts}{}
\begin{itemize}
\tightlist
\item
  Upcasts \(\leadsto\) pure functions

  \begin{itemize}
  \tightlist
  \item
    \(\sem{\upc c} : \sem{A} \to \sem{A'}\)
  \end{itemize}
\item
  Downcasts \(\leadsto\) homomorphisms of error domains

  \begin{itemize}
  \tightlist
  \item
    \(\sem{\dnc c} : \li\sem{A'} \multimap \li\sem{A}\).
  \end{itemize}
\end{itemize}

\pause

\begin{itemize}
\tightlist
\item
  Reflexivity casts \(\leadsto\) identities

  \begin{itemize}
  \tightlist
  \item
    \(\sem{\upc{r(A)}} = \id_{\sem{A}}\)
  \item
    \(\sem{\dnc{r(A)}} = \id_{\li \sem{A}}\)
  \end{itemize}
\end{itemize}

\pause

\begin{itemize}
\tightlist
\item
  Transitivity of type precision \(\leadsto\) composition

  \begin{itemize}
  \tightlist
  \item
    \(\sem{\upc{(c \comp c')}} = \sem{\upc{c'}} \circ \sem{\upc{c}}\)
  \item
    \(\sem{\dnc{(c \comp c')}} = \sem{\dnc{c}} \circ \sem{\dnc{c'}}\)
  \end{itemize}
\end{itemize}

\pause

\begin{itemize}
\tightlist
\item
  Upcasts for products/downcasts for functions \(\leadsto\) functorial
  actions of \(\times\) and \(\to\)
\end{itemize}

\pause

\begin{itemize}
\tightlist
\item
  \textbf{What about downcasts for products and upcasts for functions?}
\end{itemize}
\end{frame}

\begin{frame}{Kleisli Actions of Type Constructors}
\protect\hypertarget{kleisli-actions-of-type-constructors}{}
\end{frame}

\begin{frame}{Casts for the Dynamic Type}
\protect\hypertarget{casts-for-the-dynamic-type}{}
\begin{itemize}
\tightlist
\item
  Recall: \(D \cong \Nat + (D \times D)\, + \laterhs U(D \to \li D)\).
\item
  Upcasts are just the injections into the coproduct, except for
  \(\iarr\) where we must precompose with \(\nxt\)
\item
  Downcasts perform pattern matching and return the value if it is of
  the correct type; otherwise error

  \begin{itemize}
  \tightlist
  \item
    In the function case, the function value is only available later
  \item
    So we must insert a \(\theta\) (i.e., take an observable step) to
    gain access to the function
  \end{itemize}
\end{itemize}
\end{frame}

\begin{frame}{Extracting a Well-Behaved Semantics}
\protect\hypertarget{extracting-a-well-behaved-semantics}{}
\begin{itemize}
\tightlist
\item
  Goal: define a partial ``big-step semantics'\,' function
\item
  \textbf{Clock quantification}: method of escaping the guarded world
  and expressing ordinary set-theoretic definitions
\item
  Idea: All guarded constructs are indexed by a clock \(k\)
\end{itemize}

\pause

\begin{itemize}
\tightlist
\item
  Recall the free error domain:
  \(U\li^k X \cong X + 1 + \later^k U(\li^k X)\)
\item
  We define a global version of the free error domain:

  \begin{itemize}
  \tightlist
  \item
    \((\li^{gl} X) := \forall k . (\li^k X)\) (ignoring the \(U\))
  \end{itemize}
\end{itemize}
\end{frame}

\hypertarget{towards-a-relational-model}{%
\section{Towards a Relational Model}\label{towards-a-relational-model}}

\begin{frame}{The Goal}
\protect\hypertarget{the-goal}{}
Provide a compositional interpretation of \emph{type} and \emph{term
precision} such that we can extract the graduality relation from our
model.
\end{frame}

\begin{frame}{A Naive Approach}
\protect\hypertarget{a-naive-approach}{}
\begin{itemize}
\tightlist
\item
  Idea: enhance our value and computation types to carry a poset
  structure
\item
  For most type formers the structure is deifned functorially
\item
  \textbf{Key design choice}: ordering on the free error domain
  \(\li A\).
\item
  First attempt: \emph{step-insensitive error ordering}
\end{itemize}

\pause

\begin{figure}
      \begin{align*}
        \mho \semltbad l &\text{ iff } \top \\
        %
        \eta\, x \semltbad \eta\, y &\text{ iff } 
            x \ltdyn y \\
        %
        \theta\, \tilde{l} \semltbad \theta\, \tilde{l'} &\text{ iff } 
            \later_t (\tilde{l}_t \semltbad \tilde{l'}_t) \\
        %
        \theta\, \tilde{l} \semltbad \mho &\text{ iff } \exists n. \theta\, \tilde{l} = \delta^n(\mho) \\
        %
        \theta\, \tilde{l} \semltbad \eta\, y &\text{ iff } \exists n. \exists x \ltdyn y.
            (\theta\, \tilde{l} = \delta^n(\eta\, x)) \\
        %
        \eta\, x \semltbad \theta\, \tilde{l'} &\text { iff }
            \exists n. \exists y \gtdyn x. (\theta\, \tilde{l'} = \delta^n (\eta\, y))
    \end{align*}
    \caption{Step-insensitive error ordering}
    \label{fig:step-insensitive-error-ordering}
\end{figure}
\end{frame}

\begin{frame}{Problem: Lack of Transitivity}
\protect\hypertarget{problem-lack-of-transitivity}{}
\begin{itemize}
\tightlist
\item
  This relation is \emph{not} a partial ordering, because it is
  \textbf{not transitive}!
\item
  Follows from the following theorem:
\end{itemize}

\pause

\begin{theorem}[No-go Theorem]\label{thm:no-go}
  Let $R$ be a binary relation on the free error domain $U(\li A)$. Suppose
  $R$ satisfies the following properties:
  \begin{enumerate}
  \item Transitivity
  \item $\theta$-congruence: If $\later_t (\tilde{x}_t \binrel{R} \tilde{y}_t)$, then $\theta(\tilde{x}) \binrel{R} \theta(\tilde{y})$.
  \item Right step-insensitivity: If $x \binrel{R} y$ then $x \binrel{R} \delta y$.
  \end{enumerate}
  Then for any $l : U(\li A)$, we have $l \binrel{R} \Omega$ where $\Omega = \fix\, \theta$.
  If $R$ is left step-insensitive instead then $\Omega \binrel{R} x$.
\end{theorem}

\pause

\begin{proof}
  By L\"ob induction: we assume that $l$ is related to $\Omega$ later, which implies that
  $\theta (\nxt\, l) \binrel{R} \theta (\nxt\, \Omega)$. Then observe that

  $$l \binrel{R} \theta (\nxt\, l) \binrel{R} \theta (\nxt\, \Omega) = \Omega.$$
\end{proof}
\end{frame}

\begin{frame}{Problem: Lack of Transitivity, Continued}
\protect\hypertarget{problem-lack-of-transitivity-continued}{}
\begin{corollary}
  Let $R$ be a binary relation on $U(\li A)$. If $R$ satisfies
  transitivity, $\theta$-congruence and left and right
  step-insensitivity, then $R$ is the total relation: $\forall x, y. x
  \binrel{R} y$.
\end{corollary}

\pause

\begin{proof}
  $x \binrel{R} \Omega$ and $\Omega \binrel{R} y$ by the previous
  theorem. Then by transitivity $x \binrel R y$.
\end{proof}

\pause

\begin{itemize}
\item
  \(\semltbad\) satisfies \(\theta\)-congruence and left and right
  step-insensitivity, so if it is also transitivie then it is the total
  relation.
\item
  But e.g., for the flat poset strucutre on \(\mathbb N\),
  \(\eta 0 \semltbad \eta 1\) is false.
\item
  Thus \(\semltbad\) is \textbf{not transitive}
\end{itemize}
\end{frame}

\begin{frame}{Why Transitivity is Necessary}
\protect\hypertarget{why-transitivity-is-necessary}{}
\begin{itemize}
\tightlist
\item
  Q: How to prove graduality compositionally?
\end{itemize}

\pause

\begin{itemize}
\tightlist
\item
  Let's try to adapt the New-Licata approach to the guarded setting.
\end{itemize}
\end{frame}

\begin{frame}{Adapting the Graduality Proof}
\protect\hypertarget{adapting-the-graduality-proof}{}
\begin{itemize}
\tightlist
\item
  Interpret types as posets
\item
  Interpret a type precision derivation \(c : A \rel A'\) as a
  \emph{poset relation} between the posets \(\sem{A}\) and \(\sem{A'}\)
\end{itemize}

\pause

\begin{itemize}
\tightlist
\item
  \emph{Downward closed}: if \(x' \ltdyn_{A} x\) and
  \(x \mathrel{c} y\), then \(x' \mathrel{c} y\).
\item
  \emph{Upward closed}: if \(x \mathrel{c} y\) and \(y \ltdyn_{A'} y'\),
  then \(x \mathrel{c} y'\).
\end{itemize}

\pause

\begin{itemize}
\tightlist
\item
  Given a poset \(A\), there is an identity poset relation \(r(A)\)
  given by the poset ordering \(\ltdyn_A\)

  \begin{itemize}
  \tightlist
  \item
    \textbf{It is a poset relation because it's transitive}
  \end{itemize}
\end{itemize}

\pause

\begin{itemize}
\tightlist
\item
  Given two poset relations \(c : A_1 \rel A_2\) and
  \(c' : A_2 \rel A_3\), we define their composition \(cc'\) in the
  usual manner
\end{itemize}
\end{frame}

\begin{frame}{Modeling Term Precision}
\protect\hypertarget{modeling-term-precision}{}
\begin{itemize}
\tightlist
\item
  For closed terms: \(\sem{M \ltdyn M' : c}\) is defined to be the
  relation \(\sem{M} \binrel{(U\li\sem{c})} \sem{M'}\)
\end{itemize}

\pause

\begin{itemize}
\tightlist
\item
  For open terms: need a relationship between the denoted monotone
  functions

  \begin{itemize}
  \tightlist
  \item
    We say that \(f \ltsq{c_i}{c_o} g\) if for all \(x : A_i\) and
    \(y : A_i'\) with \(x \binrel{c_i} y\), we have
    \(f(x) \binrel{c_o} g(y)\)
  \item
    We call this relationship a \emph{square}
  \end{itemize}
\end{itemize}

\pause

\[
\begin{tikzcd}[ampersand replacement=\&]
  {A_i} \& {A_i'} \\
  {A_o} \& {A_o'}
  \arrow["{c_i}", "\shortmid"{marking}, no head, from=1-1, to=1-2]
  \arrow["f"', from=1-1, to=2-1]
  \arrow["g", from=1-2, to=2-2]
  \arrow["{c_o}"', "\shortmid"{marking}, no head, from=2-1, to=2-2]
\end{tikzcd}
\]
\end{frame}

\begin{frame}{Squares}
\protect\hypertarget{squares}{}
\begin{itemize}
\tightlist
\item
  ``vertical identity'\,' square \(\id \ltsq{c}{c} \id\) for any \(c\)
\item
  ``horizontal identity'\,' square \(f \ltsq{r(A_i)}{r(A_o)} f\) for any
  monotone function \(f : A_i \to A_o\)
\end{itemize}

\vspace{4ex}

\pause

\begin{columns}[T]
\begin{column}{0.48\textwidth}
\textbf{Vertical composition}:

If \(f \ltsq{c_1}{c_2} f'\) and \(g \ltsq{c_2}{c_3} g'\), then
\(g \circ f \ltsq{c_1}{c_3} g' \circ f'\)

% https://q.uiver.app/#q=WzAsNixbMCwwLCJBXzEiXSxbMCwxLCJBXzIiXSxbMCwyLCJBXzMiXSxbMSwwLCJBXzEnIl0sWzEsMSwiQV8yJyJdLFsxLDIsIkFfMyciXSxbMCwzLCJjXzEiLDAseyJzdHlsZSI6eyJib2R5Ijp7Im5hbWUiOiJiYXJyZWQifSwiaGVhZCI6eyJuYW1lIjoibm9uZSJ9fX1dLFsxLDQsImNfMiIsMCx7InN0eWxlIjp7ImJvZHkiOnsibmFtZSI6ImJhcnJlZCJ9LCJoZWFkIjp7Im5hbWUiOiJub25lIn19fV0sWzIsNSwiY18zIiwwLHsic3R5bGUiOnsiYm9keSI6eyJuYW1lIjoiYmFycmVkIn0sImhlYWQiOnsibmFtZSI6Im5vbmUifX19XSxbMCwxLCJmIiwyXSxbMSwyLCJnIiwyXSxbMyw0LCJmJyJdLFs0LDUsImcnIl1d
\[\begin{tikzcd}[ampersand replacement=\&]
    {A_1} \& {A_1'} \\
    {A_2} \& {A_2'} \\
    {A_3} \& {A_3'}
    \arrow["{c_1}", "\shortmid"{marking}, no head, from=1-1, to=1-2]
    \arrow["f"', from=1-1, to=2-1]
    \arrow["{f'}", from=1-2, to=2-2]
    \arrow["{c_2}", "\shortmid"{marking}, no head, from=2-1, to=2-2]
    \arrow["g"', from=2-1, to=3-1]
    \arrow["{g'}", from=2-2, to=3-2]
    \arrow["{c_3}", "\shortmid"{marking}, no head, from=3-1, to=3-2]
\end{tikzcd}
\]
\end{column}

\begin{column}{0.48\textwidth}
\textbf{Horizontal composition}:

If \(f \ltsq{c_i}{c_o} g\) and \(g \ltsq{c_i'}{c_o'} h\), then
\(f \ltsq{c_i c_i'}{c_o c_o'} h\).

% https://q.uiver.app/#q=WzAsNixbMCwwLCJBX2kiXSxbMSwwLCJBX2knIl0sWzAsMSwiQV9vIl0sWzEsMSwiQV9vJyJdLFsyLDAsIkFfaScnIl0sWzIsMSwiQV9vJyciXSxbMCwxLCJjX2kiLDAseyJzdHlsZSI6eyJib2R5Ijp7Im5hbWUiOiJiYXJyZWQifSwiaGVhZCI6eyJuYW1lIjoibm9uZSJ9fX1dLFsxLDQsImNfaSciLDAseyJzdHlsZSI6eyJib2R5Ijp7Im5hbWUiOiJiYXJyZWQifSwiaGVhZCI6eyJuYW1lIjoibm9uZSJ9fX1dLFsyLDMsImNfbyIsMix7InN0eWxlIjp7ImJvZHkiOnsibmFtZSI6ImJhcnJlZCJ9LCJoZWFkIjp7Im5hbWUiOiJub25lIn19fV0sWzMsNSwiY19vJyIsMix7InN0eWxlIjp7ImJvZHkiOnsibmFtZSI6ImJhcnJlZCJ9LCJoZWFkIjp7Im5hbWUiOiJub25lIn19fV0sWzAsMiwiZiIsMl0sWzEsMywiZyIsMl0sWzQsNSwiaCIsMl1d
\[\begin{tikzcd}[ampersand replacement=\&]
    {A_i} \& {A_i'} \& {A_i''} \\
    {A_o} \& {A_o'} \& {A_o''}
    \arrow["{c_i}", "\shortmid"{marking}, no head, from=1-1, to=1-2]
    \arrow["f"', from=1-1, to=2-1]
    \arrow["{c_i'}", "\shortmid"{marking}, no head, from=1-2, to=1-3]
    \arrow["g"', from=1-2, to=2-2]
    \arrow["h"', from=1-3, to=2-3]
    \arrow["{c_o}"', "\shortmid"{marking}, no head, from=2-1, to=2-2]
    \arrow["{c_o'}"', "\shortmid"{marking}, no head, from=2-2, to=2-3]
\end{tikzcd}\]
\end{column}
\end{columns}
\end{frame}

\begin{frame}{The Graduality Proof}
\protect\hypertarget{the-graduality-proof}{}
\begin{itemize}
\tightlist
\item
  Goal: interpret \textbf{term precision rules}
\item
  Most cases are just congruence rules and are easy to establish
\end{itemize}

\begin{mathpar}
  \inferrule{M \ltdyn M' : c_i \ra c_o \and N \ltdyn N' : c_i}{M\, N \ltdyn M'\, N' : c_o}
\end{mathpar}

\pause

\begin{itemize}
\tightlist
\item
  The core of the proof of the graduality: proving the validity of the
  \textbf{cast rules} \(\upl\), \(\upr\), \(\dnl\), and \(\dnr\).
\end{itemize}
\end{frame}

\begin{frame}{Representable Relations}
\protect\hypertarget{representable-relations}{}
\begin{itemize}
\tightlist
\item
  The four cast rules specify a relationship between the semantics of
  type precision derivations \(c\) and the corresponding casts
  \(\upc c, \dnc c\)
\end{itemize}

\pause

\begin{itemize}
\tightlist
\item
  E.g. \(\upl\) rule:
\end{itemize}

\[\begin{tikzcd}[ampersand replacement=\&]
  {\sem{A_1}} \& {\sem{A_2}} \& {\sem{A_3}} \\
  {\sem{A_2}} \&\& {\sem{A_3}}
  \arrow["\sem{c}", "\shortmid"{marking}, no head, from=1-1, to=1-2]
  \arrow["{\sem{\upc{c}}}"', from=1-1, to=2-1]
  \arrow["{\sem{c'}}", "\shortmid"{marking}, no head, from=1-2, to=1-3]
  \arrow["\id", from=1-3, to=2-3]
  \arrow["{\sem{c'}}"', "\shortmid"{marking}, no head, from=2-1, to=2-3]
\end{tikzcd}\]

\pause

\begin{itemize}
\tightlist
\item
  Problem: this rule quantifies over an arbitrary relation \(\sem{c'}\)!
\item
  This means the definition of relation is self-referential!
\end{itemize}
\end{frame}

\begin{frame}{Simpler Squares}
\protect\hypertarget{simpler-squares}{}
\begin{itemize}
\tightlist
\item
  In the presence of transitivity, we can \emph{derive} the above
  squares from simpler ones, shown below:
\end{itemize}

\begin{center}
  \begin{tabular}{ c | c }
    % UpL
    % https://q.uiver.app/#q=WzAsNCxbMCwwLCJBXzEiXSxbMCwxLCJBXzIiXSxbMSwwLCJBXzIiXSxbMSwxLCJBXzIiXSxbMCwyLCJjIiwwLHsic3R5bGUiOnsiYm9keSI6eyJuYW1lIjoiYmFycmVkIn0sImhlYWQiOnsibmFtZSI6Im5vbmUifX19XSxbMSwzLCJyKEFfMikiLDIseyJzdHlsZSI6eyJib2R5Ijp7Im5hbWUiOiJiYXJyZWQifSwiaGVhZCI6eyJuYW1lIjoibm9uZSJ9fX1dLFswLDEsIlxcdXBje2N9IiwyXSxbMiwzLCJcXGlkIl0sWzYsNywiXFx0ZXh0e1VwTH0iLDMseyJzaG9ydGVuIjp7InNvdXJjZSI6MjAsInRhcmdldCI6MjB9LCJzdHlsZSI6eyJib2R5Ijp7Im5hbWUiOiJub25lIn0sImhlYWQiOnsibmFtZSI6Im5vbmUifX19XV0=
    \begin{tikzcd}[ampersand replacement=\&]
      {A_1} \& {A_2} \\
      {A_2} \& {A_2}
      \arrow["c", "\shortmid"{marking}, no head, from=1-1, to=1-2]
      \arrow[""{name=0, anchor=center, inner sep=0}, "{\upc{c}}"', from=1-1, to=2-1]
      \arrow[""{name=1, anchor=center, inner sep=0}, "\id", from=1-2, to=2-2]
      \arrow["{r(A_2)}"', "\shortmid"{marking}, no head, from=2-1, to=2-2]
      \arrow["{\text{UpL}}"{marking, allow upside down}, draw=none, from=0, to=1]
  \end{tikzcd} &
    %
    % UpR
    % https://q.uiver.app/#q=WzAsNCxbMCwwLCJBXzEiXSxbMCwxLCJBXzEiXSxbMSwwLCJBXzEiXSxbMSwxLCJBXzIiXSxbMCwyLCJyKEFfMSkiLDAseyJzdHlsZSI6eyJib2R5Ijp7Im5hbWUiOiJiYXJyZWQifSwiaGVhZCI6eyJuYW1lIjoibm9uZSJ9fX1dLFsxLDMsImMiLDIseyJzdHlsZSI6eyJib2R5Ijp7Im5hbWUiOiJiYXJyZWQifSwiaGVhZCI6eyJuYW1lIjoibm9uZSJ9fX1dLFswLDEsIlxcaWQiLDJdLFsyLDMsIlxcdXBjIGMiXSxbNiw3LCJcXHRleHR7VXBSfSIsMyx7InNob3J0ZW4iOnsic291cmNlIjoyMCwidGFyZ2V0IjoyMH0sInN0eWxlIjp7ImJvZHkiOnsibmFtZSI6Im5vbmUifSwiaGVhZCI6eyJuYW1lIjoibm9uZSJ9fX1dXQ==
    \begin{tikzcd}[ampersand replacement=\&]
      {A_1} \& {A_1} \\
      {A_1} \& {A_2}
      \arrow["{r(A_1)}", "\shortmid"{marking}, no head, from=1-1, to=1-2]
      \arrow[""{name=0, anchor=center, inner sep=0}, "\id"', from=1-1, to=2-1]
      \arrow[""{name=1, anchor=center, inner sep=0}, "{\upc c}", from=1-2, to=2-2]
      \arrow["c"', "\shortmid"{marking}, no head, from=2-1, to=2-2]
      \arrow["{\text{UpR}}"{marking, allow upside down}, draw=none, from=0, to=1]
    \end{tikzcd} \\ \hline
    %
    %
    %
    % DnL
    % https://q.uiver.app/#q=WzAsNCxbMCwwLCJBXzIiXSxbMCwxLCJBXzEiXSxbMSwwLCJBXzIiXSxbMSwxLCJBXzIiXSxbMCwyLCJyKEFfMikiLDAseyJzdHlsZSI6eyJib2R5Ijp7Im5hbWUiOiJiYXJyZWQifSwiaGVhZCI6eyJuYW1lIjoibm9uZSJ9fX1dLFsxLDMsImMiLDIseyJzdHlsZSI6eyJib2R5Ijp7Im5hbWUiOiJiYXJyZWQifSwiaGVhZCI6eyJuYW1lIjoibm9uZSJ9fX1dLFswLDEsIlxcZG5jIGMiLDJdLFsyLDMsIlxcaWQiXSxbNiw3LCJcXHRleHR7RG5MfSIsMyx7InNob3J0ZW4iOnsic291cmNlIjoyMCwidGFyZ2V0IjoyMH0sInN0eWxlIjp7ImJvZHkiOnsibmFtZSI6Im5vbmUifSwiaGVhZCI6eyJuYW1lIjoibm9uZSJ9fX1dXQ==
    \begin{tikzcd}[ampersand replacement=\&]
      {A_2} \& {A_2} \\
      {A_1} \& {A_2}
      \arrow["{r(A_2)}", "\shortmid"{marking}, no head, from=1-1, to=1-2]
      \arrow[""{name=0, anchor=center, inner sep=0}, "{\dnc c}"', from=1-1, to=2-1]
      \arrow[""{name=1, anchor=center, inner sep=0}, "\id", from=1-2, to=2-2]
      \arrow["c"', "\shortmid"{marking}, no head, from=2-1, to=2-2]
      \arrow["{\text{DnL}}"{marking, allow upside down}, draw=none, from=0, to=1]
    \end{tikzcd} &
    %
    % DnR
    % https://q.uiver.app/#q=WzAsNCxbMCwwLCJBXzEiXSxbMCwxLCJBXzEiXSxbMSwwLCJBXzIiXSxbMSwxLCJBXzEiXSxbMCwyLCJjIiwwLHsic3R5bGUiOnsiYm9keSI6eyJuYW1lIjoiYmFycmVkIn0sImhlYWQiOnsibmFtZSI6Im5vbmUifX19XSxbMSwzLCJyKEFfMSkiLDIseyJzdHlsZSI6eyJib2R5Ijp7Im5hbWUiOiJiYXJyZWQifSwiaGVhZCI6eyJuYW1lIjoibm9uZSJ9fX1dLFswLDEsIlxcaWQiLDJdLFsyLDMsIlxcZG5jIGMiXSxbNiw3LCJcXHRleHR7RG5SfSIsMyx7InNob3J0ZW4iOnsic291cmNlIjoyMCwidGFyZ2V0IjoyMH0sInN0eWxlIjp7ImJvZHkiOnsibmFtZSI6Im5vbmUifSwiaGVhZCI6eyJuYW1lIjoibm9uZSJ9fX1dXQ==
    \begin{tikzcd}[ampersand replacement=\&]
      {A_1} \& {A_2} \\
      {A_1} \& {A_1}
      \arrow["c", "\shortmid"{marking}, no head, from=1-1, to=1-2]
      \arrow[""{name=0, anchor=center, inner sep=0}, "\id"', from=1-1, to=2-1]
      \arrow[""{name=1, anchor=center, inner sep=0}, "{\dnc c}", from=1-2, to=2-2]
      \arrow["{r(A_1)}"', "\shortmid"{marking}, no head, from=2-1, to=2-2]
      \arrow["{\text{DnR}}"{marking, allow upside down}, draw=none, from=0, to=1]
    \end{tikzcd}
  \end{tabular}
\end{center}

\begin{itemize}
\tightlist
\item
  The squares say that the relation \(c\) is \emph{representable} by the
  morphisms \(\upc{c}\) and \(\dnc{c}\)
\item
  i.e., the relation is a kind of \emph{graph} of the morphism in that
  \(x \binrel{\sem c} y\) if and only if
  \(\sem{\upc c}(x) \leq_{A_2} y\)
\end{itemize}
\end{frame}

\begin{frame}{Composing the Squares}
\protect\hypertarget{composing-the-squares}{}
Derive the original \(\upl\) square from the simpler \(\upl\) square:

\[\begin{tikzcd}[ampersand replacement=\&]
    {A_1} \& {A_2} \& {A_3} \\
    {A_2} \& {A_2} \& {A_3}
    \arrow["c", "\shortmid"{marking}, no head, from=1-1, to=1-2]
    \arrow["{\upc{c}}"', from=1-1, to=2-1]
    \arrow["{c'}", "\shortmid"{marking}, no head, from=1-2, to=1-3]
    \arrow["\id"', from=1-2, to=2-2]
    \arrow["\id", from=1-3, to=2-3]
    \arrow["{r(A_2)}"', "\shortmid"{marking}, no head, from=2-1, to=2-2]
    \arrow["{c'}"', "\shortmid"{marking}, no head, from=2-2, to=2-3]
\end{tikzcd}\]

The composition of \(r(A_2)\) with \(c'\) is equal to \(c'\), because
\(c'\) is \emph{downward-closed}
\end{frame}

\begin{frame}{Key Takeaway}
\protect\hypertarget{key-takeaway}{}
To carry out the graduality proof compositionally, the relations need to
be closed under the ordering on each side

\textbf{In particular, the ordering relations \(r(A)\) must be
transitive}
\end{frame}

\begin{frame}{Resolution: Splitting the Ordering}
\protect\hypertarget{resolution-splitting-the-ordering}{}
\begin{itemize}
\tightlist
\item
  Recall: The step-insensitive error ordering is not transitive
\end{itemize}
\end{frame}

\begin{frame}{Accounting for Steps Taken by Casts}
\protect\hypertarget{accounting-for-steps-taken-by-casts}{}
\end{frame}

\hypertarget{a-relational-model-of-gradual-typing}{%
\section{A Relational Model of Gradual
Typing}\label{a-relational-model-of-gradual-typing}}

\begin{frame}{Phase 1: Predomains and Error Domains}
\protect\hypertarget{phase-1-predomains-and-error-domains}{}
\end{frame}

\begin{frame}{Value and Computation Objects}
\protect\hypertarget{value-and-computation-objects}{}
\end{frame}

\begin{frame}{Quasi-Representable Relations}
\protect\hypertarget{quasi-representable-relations}{}
\end{frame}

\begin{frame}{Manipulating Perturbations: Push-Pull}
\protect\hypertarget{manipulating-perturbations-push-pull}{}
\end{frame}

\begin{frame}{Summary of the Model}
\protect\hypertarget{summary-of-the-model}{}
\end{frame}

\begin{frame}{Relational Adequacy}
\protect\hypertarget{relational-adequacy}{}
\end{frame}

\hypertarget{discussion-and-future-work}{%
\section{Discussion and Future Work}\label{discussion-and-future-work}}

\begin{frame}{Related Work}
\protect\hypertarget{related-work}{}
\begin{itemize}
\item
  Embedding-projection pairs
\item
  Denotational models using classical domain theory
\item
  Step-indexed operational models
\end{itemize}
\end{frame}

\begin{frame}{Agda Mechanization}
\protect\hypertarget{agda-mechanization}{}
\begin{itemize}
\tightlist
\item
  Formalized most of the key pieces of the model construction in Guarded
  Cubical Agda \cite{veltri-vezzosi2020}

  \begin{itemize}
  \tightlist
  \item
    Predomain/error domains, morphisms, relations, squares
  \item
    Free error domain
  \item
    Value/computation objects, syntactic perturbations, semantic
    perturbations, push-pull property
  \item
    Quasi-representable relations
  \item
    The dynamic type as a value object, along with the relations
    corresponding to the three injections
  \item
    Adequacy of the model with respect to the graduality property
  \end{itemize}
\end{itemize}

\pause

\begin{itemize}
\tightlist
\item
  Some formalization work remains:

  \begin{itemize}
  \tightlist
  \item
    Lemmas about quasi-representability of actions of \(\to\) and
    \(\times\) on relations
  \item
    Verifying type precision equivalence rules in the semantics
  \end{itemize}
\end{itemize}
\end{frame}

\begin{frame}{Future Work}
\protect\hypertarget{future-work}{}
\begin{itemize}
\tightlist
\item
  Extend to languages with higher-order store and/or dynamic type-tag
  generation
\item
  Prove graduality for alternative cast semantics, e.g., eager
  (\cite{herman-tomb-flanagan-2010}) and transient (\cite{transient})
  cast semantics
\item
  Utilize the \emph{intensional} nature of our model to analyze the
  efficiency of different cast semantics
\end{itemize}
\end{frame}

\end{document}
